\documentclass[10pt,a4paper]{amsart}
\usepackage[margin=19mm]{geometry}
\usepackage{amsmath,amssymb,mathtools}
\usepackage[scr=rsfs,cal=euler]{mathalfa}
\usepackage{xfrac}
\usepackage[unicode,hidelinks,bookmarks=false]{hyperref}
\newcommand{\ie}{i.e.\ }
\newcommand{\cf}{cf.\ }
\newcommand{\resp}{respectively\ }
\newcommand{\Subsection}{\subsection}
\providecommand{\nobreakdash}{\nobreak\hbox{-}}
\setlength{\emergencystretch}{2em}
\multlinegap0pt
\usepackage{amssymb}
\usepackage{enumerate}
\usepackage{dsfont}
\newtheorem{proposition}{Proposition}
\newtheorem{theorem}{Theorem}
\newcommand{\Rm}{\mathbb R^d}
\def \vf{\varphi}
\title{The Schr\"odinger Ornstein--Uhlenbeck flow: dispersion, restriction and nonlinear dynamics}
\author{Nicola Garofalo}
\date{Source: arXiv:2608.20577v1 (CC BY 4.0)}
\begin{document}
\maketitle

\noindent\emph{Source and licence.} The Schr\"odinger Ornstein--Uhlenbeck flow: dispersion, restriction and nonlinear dynamics. Version arXiv:2608.20577v1 (CC BY 4.0), distributed by arXiv under the Creative Commons Attribution 4.0 International licence, http://creativecommons.org/licenses/by/4.0/, as recorded in the arXiv licence metadata for this exact version and shown on the abstract page https://arxiv.org/abs/2608.20577v1. Full licence notice: \texttt{licenses/arxiv-2608.20577-CC-BY-4.0.txt}.\par\smallskip

\noindent\textit{Editorial scope.} Counted unit: the article's proposition P:main (source0.tex lines 894--901). The article's complete original proof of it is delivered with it (source0.tex lines 1773--1826, one \texttt{proof} environment of 54 lines, which the article places away from the statement). No mathematics is written, repaired or re-derived: the statement and the proof are the article's own lines, in the article's own notation, and no retained line is substituted or re-flowed.  Retained uncounted as the source-local context the proof needs: lines 788--790; lines 867--877; lines 1002--1004; lines 1762--1768.  Nothing was omitted from any retained span, so the package carries no omission record.  No retained line contains a citation, so the wrapper carries no bibliography and no bibliography entry.  No retained line contains a figure environment, an image-inclusion command or drawing code, so no image is delivered and the package declares no asset dependency.  Editorial formatting only, in addition: an AMS \texttt{amsart} preamble that restates the article's own declarations and macros used by the retained lines, namely the environments \texttt{proposition}, \texttt{theorem} on separate counters, together with the standard packages those lines need.  The retained environments print in this document as Proposition 1, Theorem 1 in order of appearance, whereas the article prints the same items with the numbering of its own full document.  The complete native source of the article as distributed in the arXiv e-print for this exact version is bundled unchanged as \texttt{sources/208187/source0.tex}.
\begin{equation}\label{psi}
\psi(x)=e^{-|x|^2/4}\vf(x).
\end{equation}

\begin{proposition}\label{P:semigroup}
Let $\vf\in\mathscr K(\Rm)$. Then for every $t\in J^+$,
\begin{equation}\label{final}
e^{-|x|^2/4}\left(e^{it\mathscr L}\vf\right)(x)
=(4\pi)^{-d/2}\frac{e^{idt/2}}{e^{i\pi d/4}(\sin t)^{d/2}}
 e^{i\cot t\,|x|^2/4}
\mathscr F\!\left(e^{i\cot t\,|\cdot|^2/4}\psi\right)
\left(\frac{x}{4\pi\sin t}\right),
\end{equation}
where $\psi$ is defined in \eqref{psi}.
\end{proposition}

\begin{equation}\label{parity-pre}
e^{i\pi\mathscr L}=\mathcal R,\qquad (\mathcal R g)(x)=g(-x),
\end{equation}

\begin{theorem}\label{T:har2}
Assume that $h:\Rm\to \mathbb C$ is a measurable function that satisfies 
\begin{equation}\label{hardyL2}
||e^{a|\cdot|^2} h||_{L^2(\Rm)} + ||e^{b|\cdot|^2} \hat h||_{L^2(\Rm)}<\infty.
\end{equation}
If $a b\ge \pi^2$, then $h\equiv 0$.  
\end{theorem} 

\begin{proposition}\label{P:main}
Assume that for some $a,b>0$,
\begin{equation}\label{L2}
\|e^{a|\cdot|^2}\vf\|_{L^2(\Rm,d\gamma)}
+\|e^{b|\cdot|^2}e^{is\mathscr L}\vf\|_{L^2(\Rm,d\gamma)}<\infty.
\end{equation}
If $ab\sin^2s\ge1/16$, then $\vf\equiv0$.
\end{proposition}

\begin{proof}[Proof of Proposition \ref{P:main}]
Let $\vf\in L^2(\Rm,d\gamma)$. By the assumption
$ab\sin^2s\geq 1/16$, one has $\sin s\neq0$. Write
\[
s=k\pi+\tau,
\qquad k\in\mathbb Z,
\qquad 0<\tau<\pi.
\]
By the group property and the parity identity \eqref{parity-pre},
\[
(e^{is\mathscr L}\vf)(x)
=(e^{i\tau\mathscr L}\vf)((-1)^k x).
\]
Since the Gaussian weights in \eqref{L2} are invariant under parity,
it therefore suffices to prove the result with $s$ replaced by
$\tau\in(0,\pi)$. With $\psi$ as in \eqref{psi}, consider
\begin{equation}\label{ht}
h_\tau(x)=e^{i\frac{\cot\tau |x|^2}{4}}\psi(x).
\end{equation}
We have
\begin{align}\label{one}
\int_{\Rm}e^{2a|x|^2}|h_\tau(x)|^2dx
&=\int_{\Rm}e^{2a|x|^2}|\psi(x)|^2dx
=\int_{\Rm}e^{2a|x|^2}|\vf(x)|^2e^{-\frac{|x|^2}{2}}dx
\\
&=(2\pi)^{\frac d2}\|e^{a|x|^2}\vf\|^2_{L^2(\Rm,d\gamma)}<\infty,
\notag
\end{align}
by \eqref{L2}. Since $\tau\in(0,\pi)$, Proposition \ref{P:semigroup}
gives
\begin{equation*}
\left|\widehat h_\tau\left(\frac{x}{4\pi\sin\tau}\right)\right|
=(4\pi)^{\frac d2}(\sin\tau)^{\frac d2}
e^{-\frac{|x|^2}{4}}|(e^{i\tau\mathscr L}\vf)(x)|.
\end{equation*}
Hence
\begin{align}\label{beauty}
&\left(\int_{\Rm}e^{2b|x|^2}
\left|\widehat h_\tau\left(\frac{x}{4\pi\sin\tau}\right)\right|^2dx\right)^{1/2}
\\
&=(2\pi)^{\frac d4}(4\pi)^{\frac d2}(\sin\tau)^{\frac d2}
\|e^{b|x|^2}e^{i\tau\mathscr L}\vf\|_{L^2(\Rm,d\gamma)}<\infty.
\notag
\end{align}
After the change of variables $x=4\pi\sin\tau\,y$,
\begin{equation}\label{nik}
\int_{\Rm}e^{2(16b\pi^2\sin^2\tau)|y|^2}
|\widehat h_\tau(y)|^2dy<\infty.
\end{equation}
Applying Theorem \ref{T:har2} to $h_\tau$ and using
$ab\sin^2\tau=ab\sin^2s\geq1/16$, we conclude that
$h_\tau\equiv0$. Hence $\psi\equiv0$ by \eqref{ht}, and therefore
$\vf\equiv0$.
\end{proof}

\end{document}
